\documentclass[11pt]{article}
\usepackage[T1]{fontenc}
\usepackage{amsmath,amssymb,booktabs,enumitem,geometry}
\geometry{margin=1in}

\title{An Executable Finite Containment Game Form for QCL}
\author{QCL AI-containment showcase}
\date{}

\newcommand{\Agents}{\mathit{N}}
\newcommand{\States}{\mathit{S}}
\newcommand{\Act}{\mathit{Act}}
\newcommand{\Out}{\mathsf{Out}}
\newcommand{\Eff}{\mathsf{E}}
\newcommand{\powerset}{\mathcal{P}}
\newcommand{\denote}[1]{\left[\!\left[ #1 \right]\!\right]}

\begin{document}
\maketitle

\begin{abstract}
This note specifies the finite game form used by the AI-containment
showcase. A deterministic executable transition function is the single source
of truth. Coalition outcome sets and the effectivity function are derived by
enumerating coalition strategies and all outsider responses; they are not
hand-authored. Quantified Coalition Logic (QCL) then checks properties of the
derived model. The result is relative to the selected bounded backend and is
not a security proof for a concrete cloud or AI deployment.
\end{abstract}

\section{Two backend boundaries}

The showcase retains a deterministic \texttt{synthetic} backend. Its
seven-agent \texttt{ContainmentSystem::transition} is the semantic authority
for fast tests, exhaustive game-form/effectivity regression, and the compact
mathematical model in the remainder of this note. It is deliberately not
replaced by the live implementation.

The \texttt{live} backend has a different semantic boundary. It starts a
small reference containment lab on loopback, executes every full profile, and
builds a transition table from observations:
\[
\text{ContainmentLab}\;\longrightarrow\;
\text{executed profiles}\;\longrightarrow\;
\text{ObservedTransitionTable}\;\longrightarrow\;
\text{ObservedGameForm}\;\longrightarrow\;
\texttt{derive\_effectivity}\;\longrightarrow\;\text{QCL}.
\]
The live table is a representation of observed behavior, not a second
containment algorithm. In particular, no live function maps an action profile
directly to \texttt{external\_access}, \texttt{credential\_obtained}, or
\texttt{secret\_exfiltrated}; those classifications require HTTP responses
and service event logs.

\subsection{Live reference lab}

The live QCL components are
\[
\{\texttt{agent\_a},\texttt{agent\_b},\texttt{shared\_service},
\texttt{egress\_gateway},\texttt{credential\_broker}\}.
\]
Each has two deterministic behaviors:
\begin{align*}
 \Act(\texttt{agent\_a})&=\{idle,attack\}, &
 \Act(\texttt{agent\_b})&=\{idle,attack\},\\
 \Act(\texttt{shared\_service})&=\{normal,fetch\}, &
 \Act(\texttt{egress\_gateway})&=\{enforce,bypass\},\\
 \Act(\texttt{credential\_broker})&=\{protect,expose\}.
\end{align*}
Thus live extraction executes $2^5=32$ profiles. The local
\texttt{external\_target} is environment, not a QCL agent. A separate
sandbox or safety-monitor action is unnecessary in this lab: containment is
implemented by the actual gateway process.

The services are small loopback HTTP processes. \texttt{shared\_service}
provides internal \texttt{/message} and \texttt{/messages}; only the
\texttt{shared-service-fetch} scenario enables restricted \texttt{/fetch}.
\texttt{egress\_gateway} implements actual deny/allow forwarding to the local
target. \texttt{credential\_broker} owns the in-memory fixture
\texttt{QCL\_DEMO\_SECRET\_42} and exposes \texttt{POST /credential} under its
configured policy. \texttt{external\_target} exposes \texttt{POST
/exfiltrate}, \texttt{GET /events}, and \texttt{POST /reset}.
Destination validation permits only this target; no Internet, metadata
service, host file, or real credential is reachable.

Before each profile, the runner resets all services and applies profile
configuration. Deterministic clients then issue real requests. The runner
collects responses and event logs, and records an observed outcome:
\begin{align*}
 \texttt{channel\_open}&\Leftarrow\text{message event observed},\\
 \texttt{external\_access}&\Leftarrow\text{target received a permitted request},\\
 \texttt{credential\_obtained}&\Leftarrow\text{client response contains the fixture},\\
 \texttt{secret\_exfiltrated}&\Leftarrow\text{target event log contains the fixture}.
\end{align*}
The final two classifications cannot be inferred from action names. A
repeated execution sanity check rejects a profile if its observation differs.

Conceptually, extraction stores
\[
\texttt{ObservedTransitionTable.transitions}:
  \texttt{JointActionProfile<LiveAction>}\longmapsto\texttt{StateId}.
\]
\texttt{ObservedGameForm::transition} only performs this lookup and
implements the existing \texttt{FiniteGameForm} interface. It then feeds the
shared \texttt{derive\_effectivity} algorithm; no live-specific effectivity
derivation exists.

The same formula is evaluated for both live scenarios:
\[
 \neg\langle\texttt{excludes(egress\_gateway)}\rangle
 \texttt{external\_access}.
\]
It passes for \texttt{hardened}. It fails for
\texttt{shared-service-fetch}: coalition
$\{\texttt{agent\_a},\texttt{shared\_service}\}$ can choose
\texttt{attack, fetch}, and an outsider completion produces an event at the
target without traversing the gateway. The witness contains the executed
profile and target event evidence.

This is bounded evidence about the reference lab. It excludes kernel and
network-stack correctness, hidden actions, unbounded traces, timing and race
behavior, deployment equivalence, and LLM cognition.

\section{Synthetic backend: agents, actions, and states}

Let
\[
  G=(\Agents,\States,\Act,\delta)
\]
be a finite concurrent game form. The agents are
\[
\Agents=\{a,b,ss,sb,eg,cb,sm\},
\]
standing for \texttt{agent\_a}, \texttt{agent\_b},
\texttt{shared\_service}, \texttt{sandbox},
\texttt{egress\_gateway}, \texttt{credential\_broker}, and
\texttt{safety\_monitor}. They are strategic components, not assumptions
about the cognition or intentionality of an LLM.

Each action set is binary:
\begin{align*}
 \Act(a)&=\{idle,attack\}, & \Act(b)&=\{idle,attack\},\\
 \Act(ss)&=\{isolate,relay\}, & \Act(sb)&=\{enforce,bypass\},\\
 \Act(eg)&=\{deny,allow\}, & \Act(cb)&=\{protect,expose\},\\
 \Act(sm)&=\{block,miss\}.
\end{align*}
Thus there are $2^7=128$ total joint action profiles. The state space is
\[
\begin{aligned}
\States=\{&start,contained,channel\_open,internet\_access,\\
            &credentials\_obtained,external\_compromise\}.
\end{aligned}
\]

The valuation $\pi:\States\to\powerset(Prop)$ is cumulative:
\begin{align*}
 \pi(contained)&=\{safe\},\\
 \pi(channel\_open)&=\{channel\},\\
 \pi(internet\_access)&=\{channel,internet,\\
   loss\_of\_containment\},\\
 \pi(credentials\_obtained)&=\{channel,internet,credentials,\\
   loss\_of\_containment\},\\
 \pi(external\_compromise)&=\{channel,internet,credentials,\\
   loss\_of\_containment,compromised\}.
\end{align*}
The input state is $start$. Other states are terminal outcome classes with
deterministic self-loops.

\section{Authoritative transition function}

For a total profile $\alpha\in\prod_{i\in\Agents}\Act(i)$, define
\[
  coordinated(\alpha) \iff
  \alpha(a)=attack\;\land\;\alpha(b)=attack\;\land\;\alpha(ss)=relay.
\]
The hardened transition $\delta_h$ is the following deterministic function:
\begin{align*}
\delta_h(s,\alpha)&=s &&\text{if }s\ne start,\\
\delta_h(start,\alpha)&=contained &&\text{if }\neg coordinated(\alpha),\\
\delta_h(start,\alpha)&=channel\_open &&\text{if }coordinated(\alpha)\\
 &&&\quad\text{and }(\alpha(sb)\ne bypass\lor\alpha(eg)\ne allow),\\
\delta_h(start,\alpha)&=internet\_access &&\text{if }coordinated(\alpha),\\
 &&&\quad\alpha(sb)=bypass,\ \alpha(eg)=allow,\ \alpha(cb)\ne expose,\\
\delta_h(start,\alpha)&=credentials\_obtained &&\text{if }coordinated(\alpha),\\
 &&&\quad\alpha(sb)=bypass,\ \alpha(eg)=allow,\ \alpha(cb)=expose,\\
 &&&\quad\alpha(sm)\ne miss,\\
\delta_h(start,\alpha)&=external\_compromise &&\text{otherwise.}
\end{align*}

The \texttt{shared-service-bypass} scenario changes only the coordinated
branch:
\[
  \delta_b(start,\alpha)=external\_compromise
  \quad\text{whenever }coordinated(\alpha),
\]
and returns \texttt{contained} otherwise. This is a synthetic common-mode
failure model (external proxy, ambient credentials, and monitor bypass), not a
reconstruction of an incident.

The Rust method \texttt{ContainmentSystem::transition} implements these rules.
Simulation, effectivity derivation, witnesses, and explanations call this
method; no second set of outcome rules is permitted.

\section{Coalition strategies and effectivity}

For $C\subseteq\Agents$, a coalition strategy is a partial profile
\[
  \alpha_C\in\prod_{i\in C}\Act(i).
\]
Outsiders choose any completion. Define the possible outcome set
\[
 \Out(s,C,\alpha_C)=
 \left\{\delta(s,\alpha_C\cup\alpha_{\Agents\setminus C})\mid
 \alpha_{\Agents\setminus C}\in
 \prod_{i\in\Agents\setminus C}\Act(i)\right\}.
\]
The set is adversarial: every outsider action is considered possible. A state
set $X$ is enforceable by $C$ at $s$ when
\[
 \exists\alpha_C\quad \Out(s,C,\alpha_C)\subseteq X.
\]

The raw effectivity family is
\[
 \mathrm{RawE}_s(C)=\{\Out(s,C,\alpha_C)\mid\alpha_C\text{ is a strategy}\}.
\]
The stored effectivity value is its inclusion-minimal antichain:
\[
 \Eff_s(C)=\min_{\subseteq}\bigl(\mathrm{RawE}_s(C)\bigr).
\]
The implementation inserts outcome sets into an antichain: it discards a new
set dominated by an existing subset and removes existing supersets. It never
materializes upward closure. Consequently,
\[
 X\text{ is enforceable by }C\text{ at }s
 \iff \exists Y\in\Eff_s(C):Y\subseteq X.
\]
Each surviving minimal set retains a witness $(C,\alpha_C,Y)$ for output.

\section{QCL interpretation}

Let $\denote{\varphi}_M=\{t\in\States\mid M,t\models\varphi\}$.
For a coalition predicate $P$:
\begin{align*}
 M,s\models\langle P\rangle\varphi
 &\iff \exists C\subseteq\Agents.
 C\models P\ \land\ \\
 &\qquad\exists Y\in\Eff_s(C):Y\subseteq
 \denote{\varphi}_M,\\[3pt]
 M,s\models[P]\varphi
 &\iff \forall C\subseteq\Agents.
 C\models P\ \Longrightarrow\ \\
 &\qquad\exists Y\in\Eff_s(C):Y\subseteq
 \denote{\varphi}_M.
\end{align*}
The box modality is primitive and evaluated directly. It is not replaced by
the Boolean dual of the diamond: the intended universal quantifier ranges
over coalitions, while the ability test still contains an existential
coalition strategy and universal outsider responses.

\section{Witnesses and counterexamples}

Failure explanations must follow the outer QCL connective. If
$\neg\langle P\rangle\varphi$ is false, its counterexample is a positive
existential witness:
\[
 \exists C\subseteq\Agents:\ C\models P\ \land\
 \exists\alpha_C:\ \Out(s,C,\alpha_C)\subseteq\denote{\varphi}_M.
\]
It prints $C$, $\alpha_C$, and every possible outcome, all of which satisfy
$\varphi$. This is not a failure of a defensive strategy.

Conversely, if $[P]\varphi$ is false, its counterexample is a coalition
$C$ satisfying $P$ that has \emph{no} enforcing strategy:
\[
 \exists C\subseteq\Agents:\ C\models P\ \land\
 \forall\alpha_C\ \exists\alpha_{\Agents\setminus C}:\
 \delta(s,\alpha_C\cup\alpha_{\Agents\setminus C})
 \notin\denote{\varphi}_M.
\]
The renderer checks this quantifier order before displaying a coalition. It
shows one coalition strategy, one adversarial outsider completion that breaks
the target, and states explicitly that every available coalition strategy has
such a completion. An offensive strategy reaching a bad state alone is never
evidence that a coalition lacks a different defensive strategy.

For example, in the bypass scenario $\{sm\}$ matches
$includes(sm)$ but cannot enforce $\neg compromised$: both \texttt{block}
and \texttt{miss} have outsider completions with $a=b=attack$ and
$ss=relay$, yielding $external\_compromise$.

\section{Teaching witness}

In the hardened model use
\[
 C=\{a,b,ss\},\qquad
 \alpha_C(a)=attack,\ \alpha_C(b)=attack,\ \alpha_C(ss)=relay.
\]
Varying the four outsider components yields the outcome set
\[
 \{channel\_open,internet\_access,credentials\_obtained,
 external\_compromise\}.
\]
Every member satisfies $channel$, so this coalition can enforce $channel$.
The outcome $channel\_open$ does not satisfy $internet$, so the same coalition
cannot enforce $internet$. This separates a single possible trace from an
adversarial guarantee.

\section{Exhaustive equivalence test}

For every $s\in\States$, $C\subseteq\Agents$, and $X\subseteq\States$, the
showcase compares the derived query with an independent brute-force predicate:
\[
 \mathrm{brute}(s,C,X)\iff
 \exists\alpha_C\;\forall\alpha_{\Agents\setminus C}:
 \delta(s,\alpha_C\cup\alpha_{\Agents\setminus C})\in X.
\]
The assertion is
\[
 \mathrm{derived.can\_enforce}(s,C,X)=\mathrm{brute}(s,C,X).
\]
The right-hand side calls $\delta$ directly and never calls $\Eff$. Since there
are six states, all $2^6=64$ target sets are enumerated. This is a finite,
exhaustive regression of the executable semantics and its effectivity summary.

\section{Scope, cost, and abstraction boundary}

The model is one-step and bounded. QCL here is not a temporal logic; terminal
self-loops do not assert a temporal safety property. Timing, retries, races,
hidden channels, cryptographic identity, Kubernetes behavior, and LLM
intentionality are outside the model.

The result proves only properties of $G$. Applying it to real infrastructure
requires a separate refinement or simulation argument connecting concrete
steps to $\delta$. Explicit enumeration costs exponentially in the number of
agents: $2^{|\Agents|}$ coalitions and, in the worst case,
$\prod_i|\Act(i)|$ full profiles. Symbolic transition relations, SAT/SMT,
BDDs, symmetry reduction, and compositional abstractions are possible future
extensions while preserving $\delta$ as semantic authority.

\end{document}
